\chapter{General Context}%
\label{chap:chapter_one}
\vspace{25pt} {
  \hrulefill \\%
  This chapter presents some various aspects about the project, our work method and the host company.\\
  The company I'll be partaking in this internship is publicly known as "Nethermind" and legally named "Demerzel Solution Limited", it is
  a company specialized in the development of blockchain based solutions, it tackles many aspect of the technology such as Tooling, development of the ecosystem and  
  Infrastructure. \\
  For the subject we will be working within the aspect of "Development of Ethereum Ecosystem", by building a Metaverse solution that 
  tackles the problems mentioned above while building on top of Ethereum blockchain.
  During this internship we will be developing a Metaverse, which is a complex program of real-time systems, and networking solutions
  that need to be reliable, fast and efficient, and most of all scalable.
  \hrulefill\\
}
\pagebreak
\section{\textnormal{Presentation of the Host company}}
\label{chap:section_one}

\subsection{The host Company}
  The company Demerzel Solutions LTD, better known as Nethermind, is a software company located in London, United Knigdom, it’s a part of the Ethereum softwares ecosystem. 

  \begin{figure}[H]%
    \center%
    
\includegraphics[width=0.45\textwidth]{Nethermind.png}%
    \caption[This is the Logo used by Nethermind]{Nethermind Logo}\label{fig:test}%
  \end{figure}

  \subsection{Domain of Activities}
  As of today the company has worked in various fields and project :
  \begin{figure}[H]%
    \center%
    \includegraphics[height=0.55\textwidth]{NethermindActs.png}%
    \caption[Graphic of Nethermind Activities]{Nethermind-Activities}\label{fig:test}%
  \end{figure}
  The company’s main focus is providing the Tooling for Ethereum smart contract developers, they are responsible of maintaining one of Ethereum’s main execution layer clients, and they are leading research in the ways of developing and improving the experience of Crypto field in general. 

  \subsection{Organization chart}

  The team I’ve assigned to is called «Atlantis», it is the Meta-verse research and development team, it’s comprised of 4 people: 
  \begin{description}\addtolength{\itemsep}{-0.35\baselineskip}%
    \item[\textbullet~\bfseries Project Manager] \hfill \\%
      Borys Tam.~%
    %
    \item[\textbullet~\bfseries Smart Contract Expert] \hfill \\%
      Anastazia Zeitceva (Left).~%
    %
    \item[\textbullet~\bfseries External Relations Specialist] \hfill \\%
      Julius Yls.~%
    %
    \item[\textbullet~\bfseries Interns] \hfill \\%
      Ayman BOUCHAREB.~\\%
    {\textbf{Other Interns:~}\emph{Parth Padawarth, Yohenba Kshetrimayum}}%
    %
  \end{description} 
  
  \begin{figure}[H]%
    \center%
    \includegraphics[height=0.45\textwidth]{Team-Hierarchy.jpg}%
    \caption[Chart of the hierarchy of the company]{Nethermind-Hierarchy}\label{fig:Image}%
  \end{figure}
  The team I’ve joined (Atlantis), is tasked by researching and conceptualizing solutions and projects in the meta-verse work field, my task inside of it was to experiment and develop the 3D space that will become known as the Meta-verse, by handling creating the various mechanics needed for a 3D gamified software, and the networking necessary for a multi-user experience inside said software, as well as suggesting solution and implementing the back-end server for the project. 

\section{\textnormal{Presentation of the Project}}
\label{chap:section_one}

  \subsection{The scope of the project}
  In order to provide an interactive life like experience in communication, we must first have a communication software, and a virtual space to hold these meetings, combined with the fact that Nethermind is a remote only company, we conclude that a virtual space is needed for the company to help strengthen relations between employees, as well as providing a monetisable service for external clients. 
  My task in the project was to develop and maintain an POC of the software, and participating in problem solving and brain storming sessions. 

  \subsection{Methodology of work}
  When it comes to the method followed to achieve the project, the team opted to not follow a singular methodology, arguing that introducing such restrictions and limits often introduces the mental obligation of following them, so we used what can be described as a hybrid of Agile and Waterfall, parallelising what can be done at the same time, and do in sequences what requires that, having the project as an Internal IP provided us with such liberty and freedom,  
  
  \begin{figure}[H]%
    \includegraphics[width=1.25\textwidth]{tasks-methodology.jpeg}%
    \caption[A screenshot of the team member tasks page]{Team-member-Tasks-Notion-Page}\label{fig:test}%
  \end{figure}
  we used the app called Notion, this application is suitable for file sharing and coordinating our team work, and within it, our project manager created a workflow emulating that of Kanban, where we have tasks created and assigned to the designated team member, and each task can be picked up and set to in progress. Every task's life cycle is as follows : \\ 
  \begin{figure}[H]%
    \includegraphics[width=1\textwidth]{task-life-cycle.png}%
    \caption[The life cycle of a task]{Task's life cycle}\label{fig:test}%
  \end{figure}
  we can model the life of a task as a state machine where the states represent, well the current iteration of the task, and the actions of the team are the input, the possible states of each task are : 
  \begin{description}\addtolength{\itemsep}{-0.35\baselineskip}%
    \item[\textbullet~\bfseries No Status] \hfill \\%
      for tasks newly created and un-assigned to any member ~%
    %
    \item[\textbullet~\bfseries Not Started] \hfill \\%
      for tasks that have been assigned to a member but not worked on %
    %
    \item[\textbullet~\bfseries In Progress] \hfill \\%
      for tasks being worked on%
    %
    \item[\textbullet~\bfseries In Review] \hfill \\%
      test phase of a task%
    %
    \item[\textbullet~\bfseries Archived] \hfill \\%
      tasks long aborted or finished, get moved to archived to avoid cluttering the view in the task table%
    %
    \item[\textbullet~\bfseries Obsolete] \hfill \\%
      for tasks abandoned and aborted by the team%
    %
    \item[\textbullet~\bfseries On Hold] \hfill \\%
      for tasks planned for the long term, non-urgent tasks, and tasks unable to be achieved%
    %
  \end{description} 
    
  \subsection{Timeline of the project}

    The project’s timeline was not established in the traditional manner, since the project falls under a newly developed tech market, the metaverse market in itself is in the beginning of its roadmap, and since my arrival was mid development some aspects of the project were pre-established by then, but overall the expected time line was as follows :  \\
    the bellow table shows a simplified version of the time line projected for the project, the main portion of the project life cycle is dedicated toward the iterative process of developing, testing and enhancing, provided it is Nethermind's IP, the project's goal and target audience might shift during development depending on how the market is shaping and the competition. 
  \begin{figure}[H]%
    \center%
    \begin{center}
        \begin{tabular}{||c c c c||} 
         \hline
         Activity & Duration & Start Date & End Date \\ [0.5ex] 
         \hline\hline
         Training / On-boarding & 3 Days & 02/02/2022 & 05/02/2022 \\ 
         \hline
         Specifications & 1 Week & 07/02/2022 & 11/02/2022 \\
         \hline
         Design & 1 Week & 14/02/2022 & 18/02/2022 \\
         \hline
         Programming & 6 Months & 21/02/2022  & 31/08/2022 \\ 
         \hline
         Testing & 6 Months & 21/02/2022 & 31/08/2022 \\
         \hline
        \end{tabular}
    \end{center}
    \caption[projected Time table of the project]{Project timeline}\label{fig:test}%
  \end{figure}
  in retrospect to this time the project went according to plan, with a main demo available before the predetermined deadline, 
  the actual time table went as follows :
  % Todo (Create the actual Table)
  \begin{figure}[H]%
    \center%
    \begin{center}
        \begin{tabular}{||c c c c c||} 
         \hline
         Activity & Details & Duration & Start Date & End Date \\ [0.5ex] 
         \hline\hline
         Training / On-boarding & - & 3 Day & 02/02/2022 & 4/02/2022 \\ 
         \hline
         Specifications & - & 2 Week & 07/02/2022 & 18/02/2022 \\
         \hline
         Design & - & 2 Week & 21/02/2022 & 04/03/2022 \\
         \hline
         Programming & Tooling Exploration & 1 Week & 06/03/2022  & 12/03/2022 \\ 
         \hline
         Programming & Added Avatars and dynamic Object & 1 Months & 15/03/2022  & 30/03/2022 \\ 
         \hline
         Programming & Added Networking & 1 Months & 01/04/2022  & 30/04/2022 \\ 
         \hline
         Programming & Added  Communication & 1 Months & 01/05/2022  & 30/05/2022 \\ 
         \hline
         Programming & Added  File Imports & 1 Months & 01/06/2022  & 30/06/2022 \\ 
         \hline
         Test & The global All hands Test & 1 Day & 21/05/2022  & 21/05/2022 \\ 
         \hline
        \end{tabular}
    \end{center}
    \caption[Actual Time table of the project]{Actual timeline}\label{fig:test}%
  \end{figure}
  The following figure shows the project timeline, with the main portion of the project life cycle dedicated to the iterative process of developing, testing and enhancing, provided it is Nethermind's IP, the project's goal and target audience might shift during development depending on how the market is shaping and the competition.
  \begin{figure}[H]%
    \includegraphics[width=1\textwidth]{Ganttz.png}%
    \caption[The Project's Timeline Gantt Diagram]{the actual timeline of the project vs the projected timeline }\label{fig:test}%
  \end{figure}

  \section{\textnormal{Conclusion}}
  To sum up we’ve discussed so far the scope of the project, which is taking
  place within the ”Nethermind” company, specifically within their Metaverse team called Atlantis
  the project is a software web communication software with realtime 3D interactions. 
  